% Pista ana-23j2-q6 · Anàlisi
% Procedència: PAU Catalunya 2023 · Convocatòria de juny · Sèrie 5 · Qüestió 6
\documentclass[11pt,a4paper]{article}
\usepackage{pista}
\pistainfo{Catalunya 2023}{Convocatòria de juny}{Sèrie 5}

\begin{document}

\section*{\textcolor{blauPista}{Qüestió 6}}

\noindent Sigui la funció $f(x) = \dfrac{ax^{2} + x + b}{x^{2} + 1}$.

\subsection*{\textit{a)} \normalfont Calculeu els valors dels paràmetres $a$ i $b$ si sabem que la gràfica de la funció $f$ té un extrem relatiu en $x = -1$ i passa pel punt $P = \bigl(-2, \frac{13}{5}\bigr)$. \hspace*{\fill}\textnormal{[1,25~punts]}}

\pista{Cada dada de l'enunciat es tradueix en una equació. «Extrem relatiu en $x = -1$» vol dir $f'(-1) = 0$, i «passa per $P$» vol dir $f(-2) = \frac{13}{5}$.

\medskip
En primer lloc, calcula $f'(x)$ amb la regla de la derivada d'un quocient:
\[
\left(\frac{u}{v}\right)' = \frac{u' \cdot v - u \cdot v'}{v^{2}}
\]
En resoldre $f'(-1) = 0$, només cal igualar a zero el numerador. Finalment, resol el sistema. No és tan difícil com sembla, perquè la primera equació et donarà una relació molt senzilla entre $a$ i $b$.}

\subsection*{\textit{b)} \normalfont Per al cas $a = b$, calculeu i classifiqueu els extrems relatius de la funció. \hspace*{\fill}\textnormal{[1,25~punts]}}

\pista{Substitueix $a = b$ a la derivada de l'apartat \textit{a)} i simplifica el numerador: veuràs que queda molt senzill. A continuació, iguala'l a zero per trobar els punts on hi pot haver un extrem.

\medskip
Finalment, fes una taula de signes de $f'$. El denominador sempre és positiu, perquè és un quadrat. Això vol dir que el signe de $f'$ depèn només del numerador. On la funció passa de créixer a decréixer hi ha un màxim, i on passa de decréixer a créixer, un mínim.}

\end{document}
